\begin{abstract}
Multimodal medical retrieval-augmented generation (RAG) still selects evidence
passively: candidates are reranked by similarity and the top few are fed to the
generator in one shot. This conflates \emph{relevance} with
\emph{answer-utility}, offers no recourse when the retrieved pool is poor, and,
in medicine, is misled by endoscopic images that are visually similar yet
clinically distinct. We recast evidence selection as an \emph{agentic} decision
process. A vision-language agent judges the value of each image and text
candidate, then either accepts the evidence for generation or rewrites the query
and re-retrieves. The reward is the increase in the generator's probability of
the correct answer once the selected evidence is provided---computable from
multiple-choice supervision alone, with no passage-level labels. A density-aware
router allocates the text/image retrieval budget on every round and suppresses a
modality when its signal is noise. We evaluate on EndoBench under a corrected
end-to-end image--text protocol (results TBD).
\end{abstract}
